% ============ 第1章 §1.9 Littlewood--Paley 理论 (原书 pp.30-31) ============

\section{Littlewood--Paley 理论}\label{sec:1.9}

长度 $2^{N}$ 的数列 $a = \{a_{j}\}_{j=0}^{2^{N}-1}$ 的 $\ell^{p}$ 范数为
\[
\|a\|_{\ell^{p}} = \left(\sum_{j=0}^{2^{N}-1}|a_{j}|^{p}\right)^{1/p},
\]
但只要 p 不是 1, 2, 一般地 p 次幂和很难处理. 实际上, 例如若 p = 1, 则不必开根号; 若 p = 2, 则利用 Rademacher 列 $\{r_{j}\}_{j=0}^{2^{N}-1}$, 可表示为
\[
\|a\|_{\ell^{2}} = \left(\int_{0}^{1}\left|\sum_{j=0}^{2^{N}-1}r_{j}(t)a_{j}\right|^{2}\dd t\right)^{1/2}.
\]
但是, 若 $p \neq 1, 2$, 开根号的操作会出现在定义中, 也没有便利的等式. Littlewood--Paley 理论可以说是当 $1 < p < \infty$ 时为变形 $\ell^{p}$ 范数而设的理论. 定义数列 $g(a) = \{g(a)_{j}\}_{j=0}^{2^{N}-1}$ 为
\[
g(a)_{j} = \left(|E_{N}(a)_{j}|^{2} + \sum_{k=0}^{N-1}|D_{k}(a)_{j}|^{2}\right)^{1/2}.
\]

\begin{theorem}[Littlewood--Paley 理论]\label{thm:1.13}
对数列 $a = \{a_{j}\}_{j=0}^{2^{N}-1}$, 当 $1 < p < \infty$ 时, 适当取不依赖于数列长度 $2^{N}$ 及 a 的分量的常数 $C_{p} > 1$, 则成立
\[
C_{p}^{-1}\|a\|_{\ell^{p}} \leq \|g(a)\|_{\ell^{p}} \leq C_{p}\|a\|_{\ell^{p}}.
\]
\end{theorem}

\begin{proof}
证明存在适当常数 $K_{p} > 0$ 使 $\|g(a)\|_{\ell^{p}} \leq K_{p}\|a\|_{\ell^{p}}$. 由 Khintchine 不等式, 对 $j = 0, 1, \ldots, 2^{N}-1$ 成立
\[
g(a)_{j} \leq K_{p}\left\{|E_{N}(a)_{j}| + \left(\int_{0}^{1}\left|\sum_{k=0}^{N-1}r_{k}(t)D_{k}(a)_{j}\right|^{p}\dd t\right)^{1/p}\right\}.
\]
这里, 由于 $r_{k}(t) = \pm 1$, 由定理~\ref{thm:1.12}, 若将 $K_{p}$ 重新取大, 则存在不依赖于 $t \in [0,1]$ 的常数 $K_{p} > 0$, 使
\[
\left\|\sum_{k=0}^{N-1}r_{k}(t)D_{k}(a)\right\|_{\ell^{p}} \leq K_{p}\|a\|_{\ell^{p}}
\]
也成立. 将这个不等式对 t 在 [0,1] 上积分, 代入刚才的不等式所得的常数重新记为 $K_{p}$, 得结论 $\|g(a)\|_{\ell^{p}} \leq K_{p}\|a\|_{\ell^{p}}$.

将这里得到的常数记为 $K_{p}$, 这次证明 $\|a\|_{\ell^{p}} \leq K_{p'}\|g(a)\|_{\ell^{p}}$. 将 b 写为 $b = \{b_{j}\}_{j=0}^{2^{N}-1}$, 由 Hölder 不等式的等号成立条件,
\[
\|a\|_{\ell^{p}} = \sup\left\{\left|\sum_{j=0}^{2^{N}-1}a_{j}b_{j}\right| : \|b\|_{\ell^{p'}} \leq 1\right\}.
\]
由式~\eqref{eq:1.15} 与 Cauchy--Schwarz 不等式,
\begin{align*}
\|a\|_{\ell^{p}} &= \sup\left\{\left|\langle E_{N}(a), E_{N}(b)\rangle + \sum_{k=0}^{N-1}\langle D_{k}(a), D_{k}(b)\rangle\right| : \|b\|_{\ell^{p'}} \leq 1\right\} \\
&\leq \sup\left\{\sum_{j=0}^{2^{N}-1}g(a)_{j}g(b)_{j} : \|b\|_{\ell^{p'}} \leq 1\right\}.
\end{align*}
如前所述, 若 $\|b\|_{\ell^{p'}} \leq 1$, 则 $\|g(b)\|_{\ell^{p'}} \leq K_{p'}\|b\|_{\ell^{p'}} \leq K_{p'}$, 所以由 Hölder 不等式,
\[
\|a\|_{\ell^{p}} \leq \sup\{\|g(a)\|_{\ell^{p}}\|g(b)\|_{\ell^{p'}} : \|b\|_{\ell^{p'}} \leq 1\} \leq K_{p'}\|g(a)\|_{\ell^{p}}.
\]
因此可知, 作为定理结论应取的 $C_{p}$ 可取为 $C_{p} = \max(K_{p}, K_{p'})$.
\end{proof}

\begin{exercise}\label{exer:1.26}
设 $e = (1, 0, 0, \ldots, 0)$ 为初项为 1, 其余 $2^{N}-1$ 项为 0 的数列. 当 $2^{k-1} \leq j < 2^{k}$ 时, 证明 $D_{k}(e)_{j} = 2^{-k-1}$, 并利用它证明 $\|g(e)\|_{\ell^{1}} \geq \frac{N-1}{2}$.
\end{exercise}
